\section{Background and related work}
\label{sec:background}

\subsection{Prevalence studies of AI-generated contributions}
\label{sec:background:prevalence}

\RESULT{Table~1: the audit targets. One row per published prevalence study, with
columns for detection strategy, artifact types scored, sample frame, reported
prevalence, and whether in-domain validation was performed. Populate from
LITERATURE.md section A once the four 2026 studies have been read in full and their
pipelines extracted. The final column is expected to read ``none'' throughout; that
column is the paper's motivation in one glance.}

\subsection{Detecting machine-generated text}
\label{sec:background:detection}

Zero-shot detection exploits a structural asymmetry: text sampled from a language model
tends to sit at a local maximum of that model's likelihood surface, while human text
does not. \citet{mitchell2023detectgpt} formalised this as probability curvature and
estimated it by perturbing the candidate text and re-scoring, which is accurate but
computationally heavy. \citet{bao2024fastdetectgpt} replaced the perturbation step with
a conditional probability curvature computed analytically, achieving comparable accuracy
at a fraction of the cost. \citet{hans2024binoculars} took a different route, contrasting
the perplexity of the text under one model with its cross-perplexity under a second,
closely related model; the ratio normalises away the intrinsic difficulty of the text
and yields the strongest published operating points.

Two properties of this family matter here. All of it is calibrated on text of
substantial length: the quantities being estimated are sample means over tokens, and
their variance grows as texts shorten. The null hypothesis is also ``written by a
human'' rather than ``not written by a language model''. Boilerplate produced by
non-LLM automation sits in the same low-perplexity region as model output, and nothing
in the method separates the two.

Code-specific detection adapts the same intuition to a different surface form.
\RESULT{DetectCodeGPT citation and one-paragraph method summary --- verify the exact
venue, listed as unresolved in LITERATURE.md section B.}

The closest work to ours in the detection family is \citet{ghaleb2026fingerprinting},
who trains a supervised classifier on \citednum{41} engineered features spanning commit messages,
pull request structure and code characteristics, and reports \citednum{97.2\%} F1 at identifying
\emph{which} of five agents authored a pull request. The task differs from ours: it
discriminates among agents on contributions already known to be agent-authored, so it
says nothing about the agent-versus-human decision or about false positives on human
text. We include it as a supervised, in-domain reference point. It estimates what is
achievable when a classifier is trained on the distribution it is tested on, and gives a
scale against which the off-the-shelf detectors can be read.

\subsection{Detector validity and bias}
\label{sec:background:validity}

The reliability of AI-text detection has been contested since the tools appeared.
\citet{sadasivan2023reliably} establish theoretical limits, showing that as the
distributions of human and machine text converge, any detector's performance approaches
chance, and that paraphrase attacks exploit this directly.

The empirical result most relevant to our setting is \citet{liang2023biased}, who found
that detectors classify writing by non-native English speakers as machine-generated at
substantially elevated rates, because the linguistic features that detectors treat as
evidence of machine authorship --- limited lexical variety, conventional phrasing,
lower syntactic complexity --- also characterise second-language writing. Open-source
development is a setting where this bias would be both amplified and consequential:
contributors are globally distributed, most write in English as a second language, and
being flagged as an AI contributor carries a governance cost.

Validation of detectors outside their original domain remains rare. In the news domain,
\citet{ansari2025echoes} apply Binoculars, Fast-DetectGPT and GPTZero to over
\citednum{40,000} articles to estimate the growth of LLM use in journalism --- the same instrument stack
as the software prevalence literature, applied with the same absence of in-domain
validation. \RESULT{If the pre-GPT-era Washington Post counterfactual baseline is
confirmed (see LITERATURE.md section C, currently unverified), cite it here as the
direct methodological precedent for our design and state explicitly what our study adds:
the software domain, code as well as prose, and the correction step.}

\subsection{Prevalence estimation under imperfect measurement}
\label{sec:background:correction}

Estimating the true frequency of a condition from an imperfect test is a solved problem
outside software engineering. \citet{rogan1978estimating} give the standard correction:
for an apparent prevalence $p_{\text{obs}}$ measured by a test with sensitivity $s$ and
specificity $c$, the true prevalence is
\begin{equation}
  p_{\text{true}} = \frac{p_{\text{obs}} + c - 1}{s + c - 1}.
  \label{eq:rogan-gladen}
\end{equation}
The estimator makes the dependence explicit. When $s + c - 1$ is small, meaning the test
barely outperforms chance, the correction amplifies uncertainty without bound and the
corrected estimate carries no information. Applying it to published figures needs the
quantities that have not been measured for this domain, which
Section~\ref{sec:validity} supplies.

\subsection{Bots and pre-LLM automation}
\label{sec:background:bots}

Machine-generated text in open-source repositories predates large language models by a
decade. Dependency managers, release automation, continuous integration reporters and
issue templates all produce text at scale, and identifying them is an established
research problem \RESULT{BoDeGHa citation, verify}. The confounder is measurable: in
AIDev, \AidevCommentBotShare{} of comments on agent-authored pull requests come from
bots. An estimate that does not separate bot text from model text measures their sum,
and a detector evaluation that does not report bot text separately cannot say which of
the two it responds to.

\subsection{Provenance and attribution infrastructure}
\label{sec:background:provenance}

\RESULT{PROV-O \citep{lebo2013provo} as the base vocabulary; SLSA and in-toto as the
industrial supply-chain counterpart; existing software-provenance modelling work. Frame
the gap: supply-chain provenance answers ``what produced this artifact'' for build
systems, but has no vocabulary for ``who or what authored this content, and on what
evidence''. Draft once the SLSA/in-toto positioning search in LITERATURE.md section F
is complete.}
